\section{Experimental Setup}
\label{sec:experiments}

% ------------------------------------------------------------------
\subsection{Environments}
\label{sec:experiments:envs}

We evaluate on two continuous-control locomotion tasks from the Gymnasium
MuJoCo suite~\citep{towers2024gymnasium}:
\begin{itemize}[nosep,leftmargin=*]
  \item \textbf{HalfCheetah-v5} ($d_s{=}17$, $d_a{=}6$): a planar running
        task with dense velocity reward.
  \item \textbf{Hopper-v5} ($d_s{=}11$, $d_a{=}3$): a single-legged hopping
        task requiring balance; episodes terminate on falls.
\end{itemize}
\todo{Optionally add Ant-v5 or AntMaze results.}

% ------------------------------------------------------------------
\subsection{Methods Compared}
\label{sec:experiments:methods}

We compare the following PPO variants, all using the same stochastic flow
policy architecture:
\begin{enumerate}[nosep,leftmargin=*]
  \item \textbf{Holistic PPO} (\Cref{sec:method:holistic}): standard PPO with
        a single ratio computed from the full chain log-probability.
  \item \textbf{Per-Step Uniform}: per-step PPO with uniform weights
        $w_k = 1/K$.
  \item \textbf{Per-Step Learned Global}: per-step PPO with learned global
        logits $\boldsymbol{\alpha}$.
  \item \textbf{Per-Step State-Dependent}: per-step PPO with the
        state-dependent weighting network $h_\phi$.
\end{enumerate}

% ------------------------------------------------------------------
\subsection{Architecture and Hyperparameters}
\label{sec:experiments:hyperparams}

The velocity field $v_\theta$ is parameterized as a 3-layer MLP with hidden
dimension~256 and SiLU activations.  The step index $k$ is embedded via a
sinusoidal positional encoding and concatenated with the state and latent
inputs.  The state-dependent weighting network $h_\phi$ is a 2-layer MLP with
hidden dimension~64.

Key hyperparameters are listed in \Cref{tab:hyperparams}.

\begin{table}[t]
  \centering
  \caption{Hyperparameters used across all experiments.}
  \label{tab:hyperparams}
  \begin{tabular}{@{}lc@{}}
    \toprule
    \textbf{Hyperparameter} & \textbf{Value} \\
    \midrule
    Flow steps $K$                  & 4 (default) \\
    Per-step noise $\sigma$         & 0.1 \\
    PPO clip $\epsilon$             & 0.2 \\
    GAE $\lambda$                   & 0.95 \\
    Discount $\gamma$               & 0.99 \\
    PPO epochs                      & 10 \\
    Minibatch size                  & 64 \\
    Rollout length                  & 2048 steps \\
    Learning rate (policy)          & $3 \times 10^{-4}$ \\
    Learning rate (value)           & $1 \times 10^{-3}$ \\
    Learning rate (weights $\alpha/\phi$) & $1 \times 10^{-3}$ \\
    Entropy coefficient             & 0.0 \\
    Value loss coefficient          & 0.5 \\
    Max gradient norm               & 0.5 \\
    Total environment steps         & $1 \times 10^{6}$ \\
    \bottomrule
  \end{tabular}
\end{table}

% ------------------------------------------------------------------
\subsection{Evaluation Protocol}
\label{sec:experiments:eval}

We train each method with 3 random seeds and evaluate every \todo{N} training
steps by running 10 deterministic episodes (using the mean velocity field
without per-step noise, i.e., $\sigma{=}0$ at evaluation time).  We report the
mean and standard deviation of the episodic return across seeds.

% ------------------------------------------------------------------
\subsection{Ablations}
\label{sec:experiments:ablations}

We conduct the following ablation studies:
\begin{itemize}[nosep,leftmargin=*]
  \item \textbf{Number of flow steps $K$}: we vary $K \in \{2, 4, 8\}$ to
        assess how the gap between holistic and per-step PPO changes with chain
        length.
  \item \textbf{Weighting mode}: we compare uniform, learned-global, and
        state-dependent weighting.
  \item \textbf{Noise scale $\sigma$}: we vary
        $\sigma \in \{0.01, 0.1, 0.5\}$ to test sensitivity to per-step
        stochasticity.
\end{itemize}
